\begin{frame}[fragile]{환경 만들기: \texttt{gym.make}}
\begin{lstlisting}
import gymnasium as gym

env = gym.make("CartPole-v1")   # 막대를 쓰러뜨리지 않고 균형 잡기
obs, info = env.reset(seed=0)   # 에피소드 시작 -> 초기 관측
print(type(env), "| 이름:", env.spec.id)
env.close()
\end{lstlisting}
  \begin{itemize}
    \item 환경은 \textbf{이름 문자열 하나}로 ---
      \texttt{EnvSpec}에 등록돼 있는 대로 인스턴스화.
    \item \texttt{-v1}, \texttt{-v2} 서명: 같은 게임이라도 물리
      파라미터$\cdot$종료 조건이 다른 \textbf{버전} 구분.
    \item \texttt{make}와 \texttt{reset} \textbf{분리}:
      \texttt{make}는 환경 \textit{객체}만 만들고 실제 시뮬레이션은
      \texttt{reset}이 \textbf{시작}.
      \vskip0.3em
      같은 환경 객체를 \textbf{여러 에피소드$\cdot$여러 시드}에 걸쳐
      다시 쓰는 것이 이 구조의 핵심.
  \end{itemize}
\end{frame}
